\documentclass[11pt]{article}

\usepackage[a4paper,margin=1in]{geometry}
\usepackage{amsmath,amssymb,amsthm,mathtools}
\usepackage{bm}
\usepackage{mathrsfs}
\usepackage{hyperref}
\usepackage{enumitem}
\usepackage{booktabs}
\usepackage{microtype}

\hypersetup{
  colorlinks=true,
  linkcolor=blue,
  citecolor=blue,
  urlcolor=blue
}

\newtheorem{theorem}{Theorem}
\newtheorem{proposition}[theorem]{Proposition}
\newtheorem{lemma}[theorem]{Lemma}
\newtheorem{corollary}[theorem]{Corollary}
\theoremstyle{definition}
\newtheorem{definition}[theorem]{Definition}
\newtheorem{assumption}[theorem]{Assumption}
\newtheorem{remark}[theorem]{Remark}

\newcommand{\R}{\mathbb{R}}
\newcommand{\E}{\mathbb{E}}
\newcommand{\Prob}{\mathbb{P}}
\newcommand{\supp}{\mathrm{supp}}
\newcommand{\norm}[1]{\left\lVert #1 \right\rVert}
\newcommand{\ip}[2]{\left\langle #1, #2 \right\rangle}
\newcommand{\mc}[1]{\mathcal{#1}}
\newcommand{\cE}{\mathcal{E}}
\newcommand{\cB}{\mathcal{B}}

\title{From Finite-Horizon Mean-Field Training to Linear Mode Connectivity\\for a Single-Gaussian Natural-Softplus Head}
\author{Working draft}
\date{\today}

\begin{document}
\maketitle

\begin{abstract}
We study the $K=1$ Gaussian distributional head for a two-layer network in one
paper-length chain from training dynamics to linear mode connectivity (LMC).
The model is practice faithful: the network predicts the first natural
coordinate $u_\Theta(x)$ directly and a raw scale field $s_\Theta(x)$, and the second natural
parameter is enforced by the softplus link
$\eta_2(x)=-\psi_{\rm sp}(s_\Theta(x))$. On the one hand, the Gaussian
negative log-likelihood is convex in the natural output
$\eta=(\eta_1,\eta_2)$, which removes the semi-convexity issues that appear in
mixture models. On the other hand, the practical softplus link means that
parameter interpolation is not the same as natural-output interpolation, so one
still needs to control a vector-head transfer gap and a softplus Jensen gap.

Our first main theorem gives a finite-horizon common-path result for noiseless
SGD: two width-$N$ runs started from the same initialization law track the same
deterministic mean-field path on every fixed time interval $[0,T]$. We also
include a complete continuous-time full-batch theorem as a fully proved
baseline. Our second main theorem turns the common-path input into a
Ferbach-style LMC theorem. After an optimal neuron permutation, the
interpolation barrier is bounded by a hidden-weight transfer term and a
softplus curvature term, both vanishing with width. We also state explicit
baseline and refined rates, including an $O(\sqrt{d/N}+1/N)$ route and, under
conditional head regularity, an $O(d/N)$ route.
\end{abstract}

\section{Introduction}

Linear mode connectivity (LMC) refers to the empirical phenomenon that two
independently trained networks often become connected by a low-loss linear path
after a suitable neuron permutation. In the two-layer mean-field setting,
Ferbach et al.~\cite{Ferbach2024} gave the first quantitative explanation of
this phenomenon for convex, Lipschitz losses. Their proof has a clean
two-stage structure:
\begin{enumerate}[label=(\roman*)]
  \item training-side control: two wide runs started from the same law remain
  close to the same deterministic mean-field trajectory;
  \item interpolation-side control: after an optimal transport alignment of the
  hidden neurons, the loss barrier along the parameter interpolation path is
  small.
\end{enumerate}

This paper develops the $K=1$ distributional analogue of that program for
heteroscedastic Gaussian regression. The output head is not scalar. The network
produces a two-coordinate raw output
\[
  z_\Theta(x)=\bigl(u_\Theta(x),s_\Theta(x)\bigr)\in\R^2,
\]
and the actual Gaussian natural parameters are
\[
  \eta_\Theta(x)=\Gamma(z_\Theta(x))
  =
  \bigl(u_\Theta(x),-\psi_{\rm sp}(s_\Theta(x))\bigr),
\]
where $\psi_{\rm sp}(s)=\lambda_{\min}+\log(1+e^s)$ keeps the second natural
coordinate negative. Here $u_\Theta$ is the first natural coordinate itself,
not the predictive mean; the latter is
\[
  \mu_\Theta(x)=\frac{u_\Theta(x)}{2\psi_{\rm sp}(s_\Theta(x))}.
\]
This is exactly the practical model one would train by ordinary gradient
descent in raw parameters.

The single-Gaussian case is the clean place to begin.
\begin{enumerate}[label=(\alph*)]
  \item There is no mixture softmax, no component permutation, and no gauge
  freedom.
  \item The Gaussian negative log-likelihood is convex in the natural
  coordinates $(\eta_1,\eta_2)$.
  \item The head is still genuinely vector valued, so the interpolation theory
  must handle a two-coordinate output map rather than the scalar-head setting
  of Ferbach et al.
  \item The practical softplus implementation breaks affine interpolation in
  natural space, so one must control the mismatch between parameter
  interpolation, raw-output interpolation, and natural-output interpolation.
\end{enumerate}

The paper therefore has two technical tasks.

\paragraph{Training side.}
We first write the raw-particle dynamics for the natural-softplus Gaussian
head, prove a finite-horizon confinement estimate, show Lipschitz regularity of
the induced McKean--Vlasov drift, and obtain a complete common-path theorem in
the continuous-time full-batch regime. This gives the exact training-side input
needed for Ferbach-style run-to-run geometry.

\paragraph{Interpolation side.}
We then prove a deterministic aligned barrier theorem. Because convexity holds
in natural-output space, the barrier is controlled by the distance from the
parameter interpolation path to the natural-output interpolation path. This
distance splits into two pieces:
\begin{enumerate}[label=(\roman*)]
  \item a \emph{vector-head transfer term} coming from hidden-weight mismatch;
  \item a \emph{softplus Jensen term} coming from the curvature of the
  reparameterization $s\mapsto -\psi_{\rm sp}(s)$.
\end{enumerate}

This yields the following theorem stack.
\begin{enumerate}[label=(T\arabic*)]
  \item A finite-horizon mean-field theorem for noiseless SGD, together with a
  complete continuous-time full-batch baseline theorem.
  \item A deterministic aligned barrier theorem for the $K=1$ vector head.
  \item A Ferbach-style LMC corollary obtained by combining the common-path
  theorem with Wasserstein-based neuron alignment.
\end{enumerate}

\paragraph{What is not in scope.}
We deliberately restrict to $K=1$ and do not use semi-convexity. The point of
this note is precisely that, in the single-Gaussian case, output convexity in
natural coordinates is exact and one can keep the argument much cleaner than in
the mixture setting.

\section{Setup}

Let $\mc{X}\subseteq\R^d$ be an input space and let $(X,Y)$ be a data pair. We
assume $|Y|\le Y_\ast$ almost surely. For notational simplicity we also take
the hidden parameter to be $\xi\in\R^d$; this is the standard case without an
explicit bias, and bias terms can be folded into $d$ by the usual augmentation.
Fix a bounded feature map $\varphi:\R^d\times\mc{X}\to\R$ satisfying
\[
  K_\varphi:=\sup_{\xi,x}|\varphi(\xi,x)|<\infty,
  \qquad
  K_{\nabla}:=\sup_{\xi,x}\norm{\nabla_\xi \varphi(\xi,x)}<\infty.
\]

\subsection{The practical single-Gaussian head}

We consider width-$N$ two-layer networks with mean-field scaling
\begin{equation}
  z_\Theta(x)
  =
  \bigl(u_\Theta(x),s_\Theta(x)\bigr)
  :=
  \frac1N\sum_{i=1}^N (a_i,c_i)\,\varphi(\xi_i,x)
  \in \R^2,
  \label{eq:raw_head}
\end{equation}
where
\[
  \Theta=((a_i,c_i,\xi_i))_{i=1}^N\in(\R\times\R\times\R^d)^N.
\]
The natural output is
\begin{equation}
  \eta_\Theta(x)
  :=
  \Gamma(z_\Theta(x))
  =
  \bigl(u_\Theta(x),-\psi_{\rm sp}(s_\Theta(x))\bigr),
  \label{eq:natural_head}
\end{equation}
with
\[
  \psi_{\rm sp}(s):=\lambda_{\min}+\log(1+e^s),
  \qquad
  \lambda_{\min}>0.
\]
Equivalently, the induced Gaussian has conditional mean and variance
\[
  \mu_\Theta(x)=\frac{u_\Theta(x)}{2\psi_{\rm sp}(s_\Theta(x))},
  \qquad
  \sigma_\Theta^2(x)=\frac{1}{2\psi_{\rm sp}(s_\Theta(x))}.
\]
Thus $u_\Theta$ is the first natural coordinate, not the predictive mean
itself.

In the corresponding mean-field description, a measure
$\rho\in\mc{P}(\R^2\times\R^d)$ induces
\begin{align}
  \eta_{1,\rho}(x)
  &:=
  \int a\,\varphi(\xi,x)\,\rho(d\theta),
  \\
  s_\rho(x)
  &:=
  \int c\,\varphi(\xi,x)\,\rho(d\theta),
  \\
  \eta_{2,\rho}(x)
  &:=
  -\psi_{\rm sp}(s_\rho(x)),
\end{align}
where $\theta=(a,c,\xi)$.

\begin{proposition}[Automatic negativity of the second natural coordinate]
\label{prop:softplus-domain-k1}
For every $\rho$ and every $x\in\mc{X}$,
\[
  \eta_{2,\rho}(x)\le -\lambda_{\min}<0.
\]
The same pointwise bound holds for every finite-width realization
\eqref{eq:natural_head}.
\end{proposition}

\begin{proof}
By definition, $\psi_{\rm sp}(s)\ge \lambda_{\min}$ for all $s\in\R$.
\end{proof}

\subsection{Gaussian natural loss}

For Gaussian natural parameters
$\eta=(\eta_1,\eta_2)\in \R\times(-\infty,0)$, define
\[
  A_{\rm G}(\eta_1,\eta_2)
  =
  -\frac{\eta_1^2}{4\eta_2}
  +
  \frac12\log\!\Bigl(\frac{-\pi}{\eta_2}\Bigr),
\]
and
\begin{equation}
  \ell_\eta(\eta;y)
  :=
  A_{\rm G}(\eta)-\eta_1 y-\eta_2 y^2.
  \label{eq:gaussian_nll_eta}
\end{equation}
The population risk is
\begin{equation}
  \cE(\Theta)
  :=
  \E\big[\ell_\eta(\eta_\Theta(X);Y)\big].
  \label{eq:population_risk}
\end{equation}

\subsection{Aligned interpolation}

Let $\Theta_A$ and $\Theta_B$ be two trained width-$N$ networks. Fix a neuron
permutation $\pi\in\mathfrak S_N$ and absorb it into network $B$, so we write
\[
  \Theta_A=((a_i^A,c_i^A,\xi_i^A))_{i=1}^N,
  \qquad
  \Theta_B=((a_i^B,c_i^B,\xi_i^B))_{i=1}^N.
\]
The parameter interpolation path is
\[
  \Theta_t
  :=
  \bigl((1-t)a_i^A+ta_i^B,\ (1-t)c_i^A+tc_i^B,\ (1-t)\xi_i^A+t\xi_i^B\bigr)_{i=1}^N.
\]

We also define
\[
  z_A:=z_{\Theta_A},
  \qquad
  z_B:=z_{\Theta_B},
  \qquad
  z_t:=z_{\Theta_t},
\]
the raw-output interpolation
\[
  \hat z_t:=(1-t)z_A+tz_B,
\]
the natural-output interpolation
\[
  \hat\eta_t:=(1-t)\eta_{\Theta_A}+t\eta_{\Theta_B},
\]
and the lifted raw interpolation
\[
  \tilde\eta_t:=\Gamma(\hat z_t).
\]

\begin{definition}[Aligned barrier]
For a fixed neuron permutation $\pi$, define
\[
  \cB(\Theta_A,\Theta_B;\pi)
  :=
  \sup_{t\in[0,1]}
  \Bigl\{
    \cE(\Theta_t)
    -
    (1-t)\cE(\Theta_A)
    -
    t\cE(\Theta_B)
  \Bigr\}.
\]
\end{definition}

\paragraph{Alignment diagnostics.}
The hidden mismatch is
\begin{equation}
  B_N
  :=
  \frac1N\sum_{i=1}^N \norm{\xi_i^A-\xi_i^B}^2.
  \label{eq:hidden_mismatch_k1}
\end{equation}
The endpoint raw-scale gap is
\begin{equation}
  \Delta_{s,N}
  :=
  \E\big[(s_{\Theta_A}(X)-s_{\Theta_B}(X))^2\big].
  \label{eq:raw_scale_gap}
\end{equation}
The endpoint head moment is
\begin{equation}
  M_V^\infty
  :=
  \max\!\Biggl\{
    \Bigl(\frac1N\sum_{i=1}^N (a_i^A)^2\Bigr)^{1/2},
    \Bigl(\frac1N\sum_{i=1}^N (c_i^A)^2\Bigr)^{1/2},
    \Bigl(\frac1N\sum_{i=1}^N (a_i^B)^2\Bigr)^{1/2},
    \Bigl(\frac1N\sum_{i=1}^N (c_i^B)^2\Bigr)^{1/2}
  \Biggr\}.
  \label{eq:head_moment_k1}
\end{equation}

\section{Output Geometry of the Single-Gaussian Head}

The comparison that matters is the standard one between mean--variance outputs
and natural outputs.

\begin{proposition}[Joint nonconvexity in mean--variance coordinates]
\label{prop:nonconvex-mu-v-k1}
For
\[
  \ell_{\mu v}(\mu,v;y)
  =
  \frac12\log v+\frac{(y-\mu)^2}{2v},
  \qquad v>0,
\]
the map $(\mu,v)\mapsto \ell_{\mu v}(\mu,v;y)$ is not convex on
$\R\times(0,\infty)$ for any fixed $y$.
\end{proposition}

\begin{proof}
Its Hessian is
\[
  \nabla^2\ell_{\mu v}(\mu,v;y)
  =
  \begin{pmatrix}
    v^{-1} & (y-\mu)v^{-2} \\
    (y-\mu)v^{-2} & -\frac{1}{2}v^{-2}+(y-\mu)^2v^{-3}
  \end{pmatrix},
\]
whose determinant is $-\frac{1}{2v^3}<0$.
\end{proof}

\begin{proposition}[Convexity of the Gaussian natural loss]
\label{prop:gaussian-natural-convexity}
For every fixed $y\in\R$, the map
\[
  \eta\mapsto \ell_\eta(\eta;y)=A_{\rm G}(\eta)-\eta_1 y-\eta_2 y^2
\]
is convex on the half-space $\{\eta_2<0\}$.
\end{proposition}

\begin{proof}
The Gaussian log-partition $A_{\rm G}$ is convex on $\{\eta_2<0\}$ because it
is the log-partition function of a regular minimal exponential family; the
linear term preserves convexity. See, for example,
Wainwright and Jordan~\cite{WainwrightJordan2008}.
\end{proof}

\begin{proposition}[Gaussian strip smoothness]
\label{prop:gaussian-smoothness-k1}
Fix $M>0$ and $\varepsilon>0$. On the strip
\[
  D_{M,\varepsilon}
  :=
  \{(\eta_1,\eta_2)\in\R^2: |\eta_1|\le M,\ \eta_2\le -\varepsilon\},
\]
the Gaussian log-partition function has globally bounded Hessian:
\[
  \sup_{(\eta_1,\eta_2)\in D_{M,\varepsilon}}
  \norm{\nabla^2 A_{\rm G}(\eta_1,\eta_2)}_{\mathrm{op}}
  \le
  L_{\rm G}(M,\varepsilon),
\]
where one may take
\[
  L_{\rm G}(M,\varepsilon)
  :=
  \frac{1}{2\varepsilon}
  +
  \frac{M}{2\varepsilon^2}
  +
  \frac{M^2}{2\varepsilon^3}
  +
  \frac{1}{2\varepsilon^2}.
\]
\end{proposition}

\begin{proof}
The Hessian entries are
\begin{align*}
  \partial_{11}^2 A_{\rm G}(\eta)
  &=
  -\frac{1}{2\eta_2},
  \\
  \partial_{12}^2 A_{\rm G}(\eta)
  &=
  \frac{\eta_1}{2\eta_2^2},
  \\
  \partial_{22}^2 A_{\rm G}(\eta)
  &=
  -\frac{\eta_1^2}{2\eta_2^3}+\frac{1}{2\eta_2^2}.
\end{align*}
Since $\eta_2\le -\varepsilon$ and $|\eta_1|\le M$, each entry is bounded in
absolute value by the corresponding term in $L_{\rm G}(M,\varepsilon)$.
\end{proof}

\begin{proposition}[Convexity and Lipschitz control on a natural strip]
\label{prop:convex-lipschitz-strip-k1}
Fix $M>0$ and $\lambda_{\min}>0$, and define
\[
  D_{M,\lambda_{\min}}
  :=
  \{(\eta_1,\eta_2)\in\R^2: |\eta_1|\le M,\ \eta_2\le -\lambda_{\min}\}.
\]
Then for every $|y|\le Y_\ast$:
\begin{enumerate}[label=(\roman*)]
  \item $\ell_\eta(\cdot;y)$ is convex on $D_{M,\lambda_{\min}}$;
  \item $\ell_\eta(\cdot;y)$ is
  $C_\eta(M,\lambda_{\min},Y_\ast)$-Lipschitz in the $\ell_\infty$ norm on
  $D_{M,\lambda_{\min}}$, where
  \begin{equation}
    C_\eta(M,\lambda_{\min},Y_\ast)
    :=
    \frac{M}{2\lambda_{\min}} + Y_\ast
    + \frac{M^2}{4\lambda_{\min}^2}
    + \frac{1}{2\lambda_{\min}}
    + Y_\ast^2.
    \label{eq:Ceta-k1}
  \end{equation}
\end{enumerate}
\end{proposition}

\begin{proof}
Convexity follows from Proposition~\ref{prop:gaussian-natural-convexity}.
For the Lipschitz bound, write
\[
  q_1(\eta;y)
  :=
  \partial_{\eta_1}\ell_\eta(\eta;y)
  =
  -\frac{\eta_1}{2\eta_2}-y
\]
and
\[
  q_2(\eta;y)
  :=
  \partial_{\eta_2}\ell_\eta(\eta;y)
  =
  \frac{\eta_1^2}{4\eta_2^2}
  -
  \frac{1}{2\eta_2}
  -
  y^2.
\]
On $D_{M,\lambda_{\min}}$ and $|y|\le Y_\ast$,
\[
  |q_1(\eta;y)|
  \le
  \frac{M}{2\lambda_{\min}}+Y_\ast,
  \qquad
  |q_2(\eta;y)|
  \le
  \frac{M^2}{4\lambda_{\min}^2}
  +
  \frac{1}{2\lambda_{\min}}
  +
  Y_\ast^2.
\]
The mean-value theorem in the $\ell_\infty$ norm proves the claim.
\end{proof}

\section{Finite-Horizon Mean-Field Training}

This section proves the training-side theorem that later feeds the LMC
argument. We first state the SGD-side theorem in the form needed for the LMC
bridge, using the standard four-dynamics machinery of
Mei--Misiakiewicz--Montanari~\cite{MeiMisiakiewiczMontanari2019} as organized
by Ferbach et al.~\cite{Ferbach2024}. We then keep the continuous-time
full-batch theorem as a fully proved special case and as the main internal
verification baseline.

\subsection{A finite-horizon SGD theorem for the natural-softplus head}

\begin{assumption}[Finite-horizon SGD regime]
\label{ass:softplus-sgd-k1}
Fix a horizon $T<\infty$. Assume:
\begin{enumerate}[label=(SG\arabic*)]
  \item the step sizes have the form
  \[
    s_k=\varepsilon\,\xi(k\varepsilon),
  \]
  where $\xi:[0,\infty)\to(0,\infty)$ is bounded and Lipschitz;
  \item the data stream $Z_{k+1}=(X_{k+1},Y_{k+1})$ is i.i.d.\ with law $P$,
  and $|Y|\le Y_\ast$ almost surely;
  \item $\varphi$ is bounded, continuously differentiable in $\xi$, and
  $\nabla_\xi\varphi$ is bounded and Lipschitz in $\xi$, uniformly in $x$;
  \item the initialization law $\rho_0$ on $\R^2\times\R^d$ has compact
  support;
  \item the second natural coordinate is implemented through the softplus link
  in \eqref{eq:natural_head}.
\end{enumerate}
\end{assumption}

\begin{definition}[One-sample stochastic drift]
\label{def:sample-drift-k1}
For $z=(x,y)\in \mc{X}\times\R$, define
\[
  \hat b(\theta,\rho;z)
  :=
  \bigl(\hat b_a(\theta,\rho;z),\hat b_c(\theta,\rho;z),\hat b_\xi(\theta,\rho;z)\bigr)
\]
by
\begin{align}
  \hat b_a(\theta,\rho;z)
  &:=
  -q_1(\eta_\rho(x);y)\,\varphi(\xi,x),
  \label{eq:sample-drift-a-k1}
  \\
  \hat b_c(\theta,\rho;z)
  &:=
  \psi'_{\rm sp}(s_\rho(x))\,q_2(\eta_\rho(x);y)\,\varphi(\xi,x),
  \label{eq:sample-drift-c-k1}
  \\
  \hat b_\xi(\theta,\rho;z)
  &:=
  -
  \bigl(
    a\,q_1(\eta_\rho(x);y)
    -
    c\,\psi'_{\rm sp}(s_\rho(x))q_2(\eta_\rho(x);y)
  \bigr)\nabla_\xi\varphi(\xi,x).
  \label{eq:sample-drift-xi-k1}
\end{align}
\end{definition}

\begin{proposition}[Unbiased raw stochastic update]
\label{prop:unbiased-sample-drift-k1}
Under Assumption~\ref{ass:softplus-sgd-k1},
\[
  b(\theta,\rho)=\E_Z[\hat b(\theta,\rho;Z)],
\]
where $b$ is the deterministic raw drift defined later in
Proposition~\ref{prop:softplus-raw-drift-k1} and $Z=(X,Y)\sim P$.
Consequently the noiseless SGD iteration can be written as
\begin{equation}
  \theta_i^{k+1}
  =
  \theta_i^k
  +
  s_k\,\hat b(\theta_i^k,\hat\rho_k^N;Z_{k+1}),
  \qquad
  \hat\rho_k^N:=\frac1N\sum_{i=1}^N \delta_{\theta_i^k}.
  \label{eq:sgd-update-k1}
\end{equation}
\end{proposition}

\begin{proof}
Take the expectation of \eqref{eq:sample-drift-a-k1}--\eqref{eq:sample-drift-xi-k1}
with respect to $Z=(X,Y)\sim P$ and compare with
\eqref{eq:drift-a-k1}--\eqref{eq:drift-xi-k1}.
\end{proof}

\begin{proposition}[Time-inhomogeneous finite-horizon confinement]
\label{prop:sgd-confinement-k1}
Assume Assumption~\ref{ass:softplus-sgd-k1} and let
\[
  \xi_{\max}:=\sup_{t\in[0,T]}\xi(t)<\infty.
\]
Define $A_t^{\rm sgd}$, $C_t^{\rm sgd}$, and $R_t^{\rm sgd}$ exactly as in
\eqref{eq:A-envelope-k1}--\eqref{eq:R-envelope-k1}, but with the right-hand
sides multiplied by $\xi_{\max}$. Then every characteristic of the
time-inhomogeneous mean-field flow
\[
  \dot\theta(t)=\xi(t)\,b(\theta(t),\rho_t)
\]
started from $\supp(\rho_0)$ remains in the compact set
\[
  B_T^{\rm sgd}
  :=
  \{(a,c,\xi): |a|\le A_T^{\rm sgd},\ |c|\le C_T^{\rm sgd},\ \norm{\xi}\le R_T^{\rm sgd}\}.
\]
In particular, if
\[
  M_T^{\rm sgd}:=K_\varphi A_T^{\rm sgd},
\]
then along the deterministic path one has
\[
  |\eta_{1,\rho_t}(x)|\le M_T^{\rm sgd},
  \qquad
  \eta_{2,\rho_t}(x)\le -\lambda_{\min},
  \qquad
  \forall x\in\mc{X},\ \forall t\in[0,T].
\]
\end{proposition}

\begin{proof}
The proof is identical to Proposition~\ref{prop:softplus-confinement-k1}, with
the additional prefactor $\xi(t)$ absorbed by $\xi_{\max}$ in the comparison
argument.
\end{proof}

\begin{proposition}[Discrete-time finite-horizon confinement for SGD]
\label{prop:sgd-discrete-confinement-k1}
Assume Assumption~\ref{ass:softplus-sgd-k1}. Define deterministic sequences
$A_k^\varepsilon$, $C_k^\varepsilon$, and $R_k^\varepsilon$ by
\begin{align}
  A_{k+1}^\varepsilon
  &=
  A_k^\varepsilon
  +
  s_k K_\varphi
  \left(
    \frac{K_\varphi A_k^\varepsilon}{2\lambda_{\min}}+Y_\ast
  \right),
  \qquad
  A_0^\varepsilon=A_0,
  \label{eq:discrete-A-envelope-k1}
  \\
  C_{k+1}^\varepsilon
  &=
  C_k^\varepsilon
  +
  s_k K_\varphi
  \left(
    \frac{K_\varphi^2 (A_k^\varepsilon)^2}{4\lambda_{\min}^2}
    +
    \frac{1}{2\lambda_{\min}}
    +
    Y_\ast^2
  \right),
  \qquad
  C_0^\varepsilon=C_0,
  \label{eq:discrete-C-envelope-k1}
  \\
  R_{k+1}^\varepsilon
  &=
  R_k^\varepsilon
  +
  s_k K_{\nabla}
  \Biggl[
    A_k^\varepsilon
    \left(
      \frac{K_\varphi A_k^\varepsilon}{2\lambda_{\min}}+Y_\ast
    \right)
    \nonumber\\
  &\qquad\qquad\qquad
    +
    C_k^\varepsilon
    \left(
      \frac{K_\varphi^2 (A_k^\varepsilon)^2}{4\lambda_{\min}^2}
      +
      \frac{1}{2\lambda_{\min}}
      +
      Y_\ast^2
    \right)
  \Biggr],
  \qquad
  R_0^\varepsilon=R_0.
  \label{eq:discrete-R-envelope-k1}
\end{align}
Then every SGD iterate of \eqref{eq:sgd-update-k1} satisfies
\[
  |a_i^k|\le A_k^\varepsilon,
  \qquad
  |c_i^k|\le C_k^\varepsilon,
  \qquad
  \norm{\xi_i^k}\le R_k^\varepsilon
  \qquad
  \forall i,\ \forall k\le T/\varepsilon.
\]
Moreover, the discrete envelopes are uniformly bounded by the continuous
majorants from Proposition~\ref{prop:sgd-confinement-k1}; in particular every
SGD iterate remains in $B_T^{\rm sgd}$.
\end{proposition}

\begin{proof}
We argue by induction on $k$. Assume
\[
  |a_i^k|\le A_k^\varepsilon,
  \qquad
  |c_i^k|\le C_k^\varepsilon,
  \qquad
  \norm{\xi_i^k}\le R_k^\varepsilon
\]
for all $i$. Then
\[
  |\eta_{1,\hat\rho_k^N}(x)|
  \le
  \frac1N\sum_{i=1}^N |a_i^k|\,|\varphi(\xi_i^k,x)|
  \le
  K_\varphi A_k^\varepsilon,
\]
and Proposition~\ref{prop:softplus-domain-k1} gives
$\eta_{2,\hat\rho_k^N}(x)\le -\lambda_{\min}$. Hence
\eqref{eq:q1-k1}--\eqref{eq:q2-k1} imply the same bounds as in the proof of
Proposition~\ref{prop:softplus-confinement-k1}, now with $A_k^\varepsilon$ in
place of $A_t$. Inserting these estimates into the single-sample update
\eqref{eq:sgd-update-k1} gives
\[
  |a_i^{k+1}|
  \le
  A_{k+1}^\varepsilon,
  \qquad
  |c_i^{k+1}|
  \le
  C_{k+1}^\varepsilon,
  \qquad
  \norm{\xi_i^{k+1}}
  \le
  R_{k+1}^\varepsilon.
\]
This closes the induction.

The comparison between the discrete recursion and the continuous majorant ODEs
is standard Euler-vs-ODE monotonicity: since $s_k\le \varepsilon\xi_{\max}$ and
the right-hand sides are nondecreasing in the envelopes, one has
$A_k^\varepsilon\le A_{k\varepsilon}^{\rm sgd}$, $C_k^\varepsilon\le
C_{k\varepsilon}^{\rm sgd}$, and $R_k^\varepsilon\le R_{k\varepsilon}^{\rm sgd}$
for all $k\le T/\varepsilon$.
\end{proof}

\begin{proposition}[Uniform stochastic-drift regularity]
\label{prop:sample-drift-regularity-k1}
Assume Assumption~\ref{ass:softplus-sgd-k1}. There exists
$\widehat L_T<\infty$ such that for every sample $z=(x,y)$ in the support of
$P$, all $\theta,\bar\theta\in B_T^{\rm sgd}$, and all probability measures
$\rho,\nu$ supported on $B_T^{\rm sgd}$,
\begin{align}
  \norm{\hat b(\theta,\rho;z)}
  &\le
  \widehat L_T,
  \label{eq:sample-drift-bound-k1}
  \\
  \norm{\hat b(\theta,\rho;z)-\hat b(\bar\theta,\nu;z)}
  &\le
  \widehat L_T\bigl(\norm{\theta-\bar\theta}+W_1(\rho,\nu)\bigr).
  \label{eq:sample-drift-lipschitz-k1}
\end{align}
\end{proposition}

\begin{proof}
On $B_T^{\rm sgd}$, the strip bound from
Proposition~\ref{prop:sgd-confinement-k1} and
Proposition~\ref{prop:gaussian-smoothness-k1} imply that
$q_1(\eta_\rho(x);y)$ and $q_2(\eta_\rho(x);y)$ are uniformly bounded and
Lipschitz in $(\theta,\rho)$ for $|y|\le Y_\ast$. Since $\psi'_{\rm sp}$ is
bounded and Lipschitz, and since $\varphi$ and $\nabla_\xi\varphi$ are bounded
and Lipschitz in $\xi$, the formulas
\eqref{eq:sample-drift-a-k1}--\eqref{eq:sample-drift-xi-k1} immediately yield
\eqref{eq:sample-drift-bound-k1} and \eqref{eq:sample-drift-lipschitz-k1}.
\end{proof}

\begin{theorem}[Finite-horizon SGD mean-field approximation]
\label{thm:sgd-mean-field-k1}
Assume Assumption~\ref{ass:softplus-sgd-k1}. Consider the noiseless SGD
iterates \eqref{eq:sgd-update-k1} and let $\rho_t$ be the deterministic
measure-valued path solving
\begin{equation}
  \partial_t \rho_t
  +
  \nabla_\theta\cdot\bigl(\rho_t\,\xi(t)b(\theta,\rho_t)\bigr)
  =
  0,
  \qquad
  \rho_{t=0}=\rho_0.
  \label{eq:sgd-pde-k1}
\end{equation}
Then:
\begin{enumerate}[label=(\roman*)]
  \item the PDE \eqref{eq:sgd-pde-k1} is well posed on $[0,T]$;
  \item there exists an error term $\delta_{N,\varepsilon}(T)\to 0$ as
  $N\to\infty$ and $\varepsilon\downarrow 0$ such that
  \[
    \sup_{k\le T/\varepsilon}
    W_1\!\bigl(\hat\rho_k^N,\rho_{k\varepsilon}\bigr)
    \le
    \delta_{N,\varepsilon}(T)
  \]
  with high probability;
  \item in bounded-support regimes, one obtains the same qualitative structure
  as in the Mei/Ferbach four-dynamics argument: a Monte-Carlo term of order
  $N^{-1/2}$ together with discretization and stochastic-gradient terms of
  order $\varepsilon$ and $\sqrt{\varepsilon}$.
\end{enumerate}
\end{theorem}

\begin{proof}[Proof sketch]
Propositions~\ref{prop:unbiased-sample-drift-k1},
\ref{prop:sgd-confinement-k1}, \ref{prop:sgd-discrete-confinement-k1}, and
\ref{prop:sample-drift-regularity-k1} verify the model-specific hypotheses
needed for the bounded-support
four-dynamics comparison of Mei--Misiakiewicz--Montanari
\cite{MeiMisiakiewiczMontanari2019}. Ferbach et al.~\cite{Ferbach2024} use the
same reduction in Appendix~C.2 for noiseless regularization-free SGD. Applying
that finite-horizon machinery to the present raw stochastic drift
\eqref{eq:sgd-update-k1} yields the stated approximation theorem.
\end{proof}

\begin{corollary}[SGD two-run common path]
\label{cor:sgd-common-path-k1}
Assume Assumption~\ref{ass:softplus-sgd-k1}. Let
$\hat\rho_{A,k}^N,\hat\rho_{B,k}^N$ be two independent width-$N$ SGD runs of
\eqref{eq:sgd-update-k1}, both initialized i.i.d.\ from the same law $\rho_0$
and trained over the same horizon $[0,T]$. Then
\[
  \sup_{k\le T/\varepsilon}
  W_1\!\bigl(\hat\rho_{A,k}^N,\hat\rho_{B,k}^N\bigr)
  \le
  2\,\delta_{N,\varepsilon}(T)
\]
with high probability. In particular,
\[
  \sup_{k\le T/\varepsilon}
  W_1\!\bigl(\hat\rho_{A,k}^N,\hat\rho_{B,k}^N\bigr)\to 0
\]
in probability as $N\to\infty$ and $\varepsilon\downarrow 0$.
\end{corollary}

\begin{proof}
Apply Theorem~\ref{thm:sgd-mean-field-k1} to each run and use the triangle
inequality.
\end{proof}

\subsection{A complete full-batch baseline theorem}

\begin{assumption}[Continuous-time full-batch regime]
\label{ass:complete-full-batch-k1}
Fix a horizon $T<\infty$. Assume:
\begin{enumerate}[label=(CF\arabic*)]
  \item training is continuous-time full-batch gradient flow on the population
  risk \eqref{eq:population_risk}, with the usual mean-field scaling;
  \item $|Y|\le Y_\ast$ almost surely;
  \item $\varphi$ is bounded and $\nabla_\xi\varphi$ is bounded and Lipschitz
  in $\xi$, uniformly in $x$;
  \item the initialization law $\rho_0$ on $\R^2\times\R^d$ has compact
  support;
  \item the second natural coordinate is implemented through the softplus link
  in \eqref{eq:natural_head}.
\end{enumerate}
\end{assumption}

For later use define
\begin{align}
  q_1(\eta;y)
  &:=
  \partial_{\eta_1}\ell_\eta(\eta;y)
  =
  -\frac{\eta_1}{2\eta_2}-y,
  \label{eq:q1-k1}
  \\
  q_2(\eta;y)
  &:=
  \partial_{\eta_2}\ell_\eta(\eta;y)
  =
  \frac{\eta_1^2}{4\eta_2^2}
  -
  \frac{1}{2\eta_2}
  -
  y^2.
  \label{eq:q2-k1}
\end{align}

\begin{proposition}[Raw drift formula]
\label{prop:softplus-raw-drift-k1}
Assume Assumption~\ref{ass:complete-full-batch-k1}. For
$\theta=(a,c,\xi)\in\R^2\times\R^d$ and a probability measure $\rho$, define
\[
  \eta_\rho(x)
  :=
  \bigl(\eta_{1,\rho}(x),\eta_{2,\rho}(x)\bigr),
  \qquad
  h_\rho(x,y)
  :=
  \psi_{\rm sp}'(s_\rho(x))\,q_2(\eta_\rho(x);y).
\]
Then the mean-field-scaled full-batch gradient flow of the population risk is
the interacting particle system
\begin{equation}
  \dot\theta_i^N(t)
  =
  b\bigl(\theta_i^N(t),\rho_t^N\bigr),
  \qquad
  \rho_t^N:=\frac1N\sum_{i=1}^N \delta_{\theta_i^N(t)},
  \label{eq:particle-system-k1}
\end{equation}
with drift $b=(b_a,b_c,b_\xi)$ given by
\begin{align}
  b_a(\theta,\rho)
  &:=
  -
  \E\big[q_1(\eta_\rho(X);Y)\,\varphi(\xi,X)\big],
  \label{eq:drift-a-k1}
  \\
  b_c(\theta,\rho)
  &:=
  \E\big[h_\rho(X,Y)\,\varphi(\xi,X)\big],
  \label{eq:drift-c-k1}
  \\
  b_\xi(\theta,\rho)
  &:=
  -
  \E\Big[
    \big(a\,q_1(\eta_\rho(X);Y)-c\,h_\rho(X,Y)\big)\nabla_\xi\varphi(\xi,X)
  \Big].
  \label{eq:drift-xi-k1}
\end{align}
\end{proposition}

\begin{proof}
Let
\[
  \mc{L}_N(\theta_1,\dots,\theta_N)
  :=
  \E\big[\ell_\eta(\eta_{\rho_N}(X);Y)\big],
  \qquad
  \rho_N=\frac1N\sum_{j=1}^N \delta_{\theta_j}.
\]
Since
\[
  \partial_{\theta_i}\eta_{1,\rho_N}(x)
  =
  \frac1N\bigl(\varphi(\xi_i,x),0,a_i\nabla_\xi\varphi(\xi_i,x)\bigr)
\]
and
\[
  \partial_{\theta_i}\eta_{2,\rho_N}(x)
  =
  -\frac1N\psi_{\rm sp}'(s_{\rho_N}(x))
  \bigl(0,\varphi(\xi_i,x),c_i\nabla_\xi\varphi(\xi_i,x)\bigr),
\]
the chain rule yields
\[
  \nabla_{\theta_i}\mc{L}_N
  =
  \frac1N
  \begin{pmatrix}
    \E\big[q_1(\eta_{\rho_N}(X);Y)\varphi(\xi_i,X)\big]
    \\
    -\E\big[h_{\rho_N}(X,Y)\varphi(\xi_i,X)\big]
    \\
    \E\big[(a_iq_1(\eta_{\rho_N}(X);Y)-c_i h_{\rho_N}(X,Y))
    \nabla_\xi\varphi(\xi_i,X)\big]
  \end{pmatrix}.
\]
Multiplying by the mean-field factor $N$ in continuous-time gradient flow gives
\eqref{eq:particle-system-k1}--\eqref{eq:drift-xi-k1}.
\end{proof}

\begin{proposition}[Finite-horizon confinement]
\label{prop:softplus-confinement-k1}
Assume Assumption~\ref{ass:complete-full-batch-k1}. Let
\[
  A_0:=\sup_{(a,c,\xi)\in\supp(\rho_0)}|a|,
  \qquad
  C_0:=\sup_{(a,c,\xi)\in\supp(\rho_0)}|c|,
  \qquad
  R_0:=\sup_{(a,c,\xi)\in\supp(\rho_0)}\norm{\xi}.
\]
Define the common initialization box
\[
  B_0
  :=
  \{(a,c,\xi): |a|\le A_0,\ |c|\le C_0,\ \norm{\xi}\le R_0\}.
\]
Define envelopes on $[0,T]$ by
\begin{align}
  \dot A_t
  &=
  K_\varphi\left(\frac{K_\varphi A_t}{2\lambda_{\min}}+Y_\ast\right),
  \qquad
  A_{t=0}=A_0,
  \label{eq:A-envelope-k1}
  \\
  \dot C_t
  &=
  K_\varphi\left(
    \frac{K_\varphi^2 A_t^2}{4\lambda_{\min}^2}
    +
    \frac{1}{2\lambda_{\min}}
    +
    Y_\ast^2
  \right),
  \qquad
  C_{t=0}=C_0,
  \label{eq:C-envelope-k1}
  \\
  \dot R_t
  &=
  K_{\nabla}\left[
    A_t\left(\frac{K_\varphi A_t}{2\lambda_{\min}}+Y_\ast\right)
    +
    C_t\left(
      \frac{K_\varphi^2 A_t^2}{4\lambda_{\min}^2}
      +
      \frac{1}{2\lambda_{\min}}
      +
      Y_\ast^2
    \right)
  \right],
  \qquad
  R_{t=0}=R_0.
  \label{eq:R-envelope-k1}
\end{align}
Then $A_t$, $C_t$, and $R_t$ remain finite on $[0,T]$. Define for each
$t\in[0,T]$
\[
  B_t
  :=
  \{(a,c,\xi): |a|\le A_t,\ |c|\le C_t,\ \norm{\xi}\le R_t\},
  \qquad
  M_t:=K_\varphi A_t.
\]
Every characteristic of any mean-field solution path of \eqref{eq:mf-pde-k1}
whose initial law is supported on $B_0$, as well as every particle path of
\eqref{eq:particle-system-k1} whose initial particles lie in $B_0$, satisfies
\[
  |a(t)|\le A_t,
  \qquad
  |c(t)|\le C_t,
  \qquad
  \norm{\xi(t)}\le R_t
  \qquad
  \forall t\in[0,T].
\]
In particular, for every such mean-field solution path,
\[
  \supp(\rho_t)\subseteq B_t\subseteq B_T
  \qquad
  \forall t\in[0,T],
\]
and
\[
  |\eta_{1,\rho_t}(x)|\le M_t\le M_T
  \qquad
  \forall x\in\mc{X},\ \forall t\in[0,T].
\]
\end{proposition}

\begin{proof}
Because $\eta_{2,\rho}(x)\le -\lambda_{\min}$ and
$|\eta_{1,\rho}(x)|\le K_\varphi A_t$, \eqref{eq:q1-k1}--\eqref{eq:q2-k1}
give
\[
  |q_1(\eta_\rho(x);y)|
  \le
  \frac{K_\varphi A_t}{2\lambda_{\min}}+Y_\ast
\]
and
\[
  |q_2(\eta_\rho(x);y)|
  \le
  \frac{K_\varphi^2 A_t^2}{4\lambda_{\min}^2}
  +
  \frac{1}{2\lambda_{\min}}
  +
  Y_\ast^2.
\]
Since $0<\psi'_{\rm sp}<1$, the drift formulas imply
\[
  |\dot a(t)|\le \dot A_t,
  \qquad
  |\dot c(t)|\le \dot C_t,
  \qquad
  \norm{\dot\xi(t)}\le \dot R_t.
\]
Scalar comparison yields the claim.
\end{proof}

\begin{lemma}[Frozen-path confinement on the common envelopes]
\label{lem:frozen-confinement-k1}
Assume Assumption~\ref{ass:complete-full-batch-k1}. Let $B_0$ and $(B_t)$ be as
in Proposition~\ref{prop:softplus-confinement-k1}. Let
$m=(m_t)_{t\in[0,T]}$ be a continuous measure path such that
\[
  \supp(m_t)\subseteq B_t
  \qquad
  \forall t\in[0,T].
\]
Then every solution of
\[
  \dot\theta(t)=b(\theta(t),m_t),
  \qquad
  \theta(0)\in B_0,
\]
satisfies
\[
  |a(t)|\le A_t,
  \qquad
  |c(t)|\le C_t,
  \qquad
  \norm{\xi(t)}\le R_t
  \qquad
  \forall t\in[0,T].
\]
In particular,
\[
  X_t^m(B_0)\subseteq B_t
  \qquad
  \forall t\in[0,T],
\]
where $X_t^m$ denotes the frozen-path flow map.
\end{lemma}

\begin{proof}
Because $\supp(m_t)\subseteq B_t$, one has
$\eta_{2,m_t}(x)\le -\lambda_{\min}$ and
$|\eta_{1,m_t}(x)|\le K_\varphi A_t$. Therefore
\eqref{eq:q1-k1}--\eqref{eq:q2-k1} give
\[
  |q_1(\eta_{m_t}(x);y)|
  \le
  \frac{K_\varphi A_t}{2\lambda_{\min}}+Y_\ast
\]
and
\[
  |q_2(\eta_{m_t}(x);y)|
  \le
  \frac{K_\varphi^2 A_t^2}{4\lambda_{\min}^2}
  +
  \frac{1}{2\lambda_{\min}}
  +
  Y_\ast^2.
\]
Since $0<\psi'_{\rm sp}<1$, the drift formulas imply
\[
  |\dot a(t)|\le \dot A_t,
  \qquad
  |\dot c(t)|\le \dot C_t,
  \qquad
  \norm{\dot\xi(t)}\le \dot R_t.
\]
Scalar comparison yields the claim.
\end{proof}

\begin{proposition}[Drift regularity on the confinement set]
\label{prop:softplus-drift-regularity-k1}
Assume Assumption~\ref{ass:complete-full-batch-k1}. There exists $L_T<\infty$
such that for all $\theta,\bar\theta\in B_T$ and all probability measures
$\rho,\nu$ supported on $B_T$,
\begin{equation}
  \norm{b(\theta,\rho)-b(\bar\theta,\nu)}
  \le
  L_T\big(\norm{\theta-\bar\theta}+W_1(\rho,\nu)\big),
  \label{eq:drift-lipschitz-k1}
\end{equation}
and
\[
  \sup_{\theta\in B_T,\ \supp(\rho)\subseteq B_T}\norm{b(\theta,\rho)}
  \le
  L_T.
\]
\end{proposition}

\begin{proof}
On $B_T$, the maps
$\theta\mapsto a\varphi(\xi,x)$ and $\theta\mapsto c\varphi(\xi,x)$ are
uniformly bounded and Lipschitz in $\theta$. Hence for every $x$,
\[
  |\eta_{1,\rho}(x)-\eta_{1,\nu}(x)| + |s_\rho(x)-s_\nu(x)|
  \le
  C_T^{(1)}W_1(\rho,\nu)
\]
for some $C_T^{(1)}<\infty$. Since $\psi'_{\rm sp}$ is bounded and Lipschitz,
and Proposition~\ref{prop:gaussian-smoothness-k1} gives a uniform Lipschitz
constant for $\nabla_\eta \ell_\eta$ on $D_{M_T,\lambda_{\min}}$, both
\[
  (\rho,x,y)\mapsto q_1(\eta_\rho(x);y),
  \qquad
  (\rho,x,y)\mapsto h_\rho(x,y)
\]
are uniformly bounded and Lipschitz in $\rho$ on the confinement set. Inserting
these estimates into \eqref{eq:drift-a-k1}--\eqref{eq:drift-xi-k1} yields
\eqref{eq:drift-lipschitz-k1}. The uniform bound is immediate from the same
estimates.
\end{proof}

\begin{theorem}[Finite-horizon mean-field theorem on the common initialization box]
\label{thm:complete-full-batch-k1}
Assume Assumption~\ref{ass:complete-full-batch-k1}, and let $B_0$ and $(B_t)$
be as in Proposition~\ref{prop:softplus-confinement-k1}. Then:
\begin{enumerate}[label=(\roman*)]
  \item for every initial measure $\mu$ supported on $B_0$, there exists a
  unique measure-valued path $(\rho_t)_{t\in[0,T]}$ such that
  $\supp(\rho_t)\subseteq B_t$ for all $t\in[0,T]$ and
  \begin{equation}
    \partial_t \rho_t
    +
    \nabla_\theta\cdot\big(\rho_t\,b(\theta,\rho_t)\big)
    =
    0,
    \qquad
    \rho_{t=0}=\mu.
    \label{eq:mf-pde-k1}
  \end{equation}
  \item for any two solutions $(\rho_t)$ and $(\nu_t)$ of
  \eqref{eq:mf-pde-k1} with initial data $\bar\rho_0$ and $\bar\nu_0$
  supported on $B_0$,
  \begin{equation}
    \sup_{t\in[0,T]}W_1(\rho_t,\nu_t)
    \le
    e^{2L_TT}W_1(\bar\rho_0,\bar\nu_0);
    \label{eq:pde-stability-k1}
  \end{equation}
  \item if $\rho_t^N$ denotes the empirical measure of the interacting particle
  system \eqref{eq:particle-system-k1} with initial particles
  $\theta_i^0\in B_0$, and if $\rho_t$ denotes the solution of
  \eqref{eq:mf-pde-k1} with initial law $\rho_0$, then
  \begin{equation}
    \sup_{t\in[0,T]}W_1(\rho_t^N,\rho_t)
    \le
    e^{2L_TT}W_1(\rho_0^N,\rho_0).
    \label{eq:particle-tracking-k1}
  \end{equation}
\end{enumerate}
Consequently, if $\theta_i^0$ are i.i.d.\ with law $\rho_0$, then
\[
  \sup_{t\in[0,T]}W_1(\rho_t^N,\rho_t)\to 0
\]
almost surely, hence also in probability, as $N\to\infty$.
\end{theorem}

\begin{proof}
Define
\[
  \mc{A}_T
  :=
  \Bigl\{
    m\in C([0,T];\mc{P}(B_T)):
    \supp(m_t)\subseteq B_t\ \forall t\in[0,T]
  \Bigr\}.
\]
Fix an initial measure $\mu$ supported on $B_0$. For a path
$m=(m_t)_{t\in[0,T]}\in\mc{A}_T$, let $X_t^m(\theta)$ solve
\[
  \frac{d}{dt}X_t^m(\theta)=b(X_t^m(\theta),m_t),
  \qquad
  X_0^m(\theta)=\theta.
\]
Lemma~\ref{lem:frozen-confinement-k1} implies that this ODE is globally defined
on $[0,T]$ for every $\theta\in B_0$, and that $X_t^m(B_0)\subseteq B_t$ for
all $t$. Define
\[
  \Gamma_\mu(m)_t:=(X_t^m)_\#\mu.
\]
Thus $\Gamma_\mu$ maps $\mc{A}_T$ into itself. If $m,n\in\mc{A}_T$, then for
every $\theta\in B_0$ Proposition~\ref{prop:softplus-drift-regularity-k1}
gives
\[
  \frac{d}{dt}\norm{X_t^m(\theta)-X_t^n(\theta)}
  \le
  L_T\norm{X_t^m(\theta)-X_t^n(\theta)}
  +
  L_TW_1(m_t,n_t).
\]
Gronwall yields
\[
  W_1(\Gamma_\mu(m)_t,\Gamma_\mu(n)_t)
  \le
  L_Te^{L_Tt}\int_0^t W_1(m_s,n_s)\,ds.
\]
Hence $\Gamma_\mu$ is a contraction on sufficiently short intervals; iterating
the fixed-point argument proves (i).

For (ii), Proposition~\ref{prop:softplus-confinement-k1} gives
$\supp(\rho_t),\supp(\nu_t)\subseteq B_t\subseteq B_T$ because the initial
measures are supported on $B_0$. Couple the two solutions through
characteristics
$X_t,Y_t$ satisfying
\[
  \dot X_t=b(X_t,\rho_t),
  \qquad
  \dot Y_t=b(Y_t,\nu_t).
\]
Then
\[
  \frac{d}{dt}\E\norm{X_t-Y_t}
  \le
  L_T\E\norm{X_t-Y_t}+L_TW_1(\rho_t,\nu_t)
  \le
  2L_T\E\norm{X_t-Y_t}.
\]
Taking the infimum over all couplings of $(X_0,Y_0)$ proves
\eqref{eq:pde-stability-k1}.

For (iii), Proposition~\ref{prop:softplus-confinement-k1} gives
$\supp(\rho_t^N)\subseteq B_t$ because the initial particles lie in $B_0$.
Moreover, $\rho_t^N$ solves \eqref{eq:mf-pde-k1} in the weak sense because,
for every smooth test function $\chi$,
\[
  \frac{d}{dt}\int \chi(\theta)\,\rho_t^N(d\theta)
  =
  \frac1N\sum_{i=1}^N
  \nabla\chi(\theta_i^N(t))\cdot b(\theta_i^N(t),\rho_t^N)
  =
  \int \nabla\chi(\theta)\cdot b(\theta,\rho_t^N)\,\rho_t^N(d\theta).
\]
Applying \eqref{eq:pde-stability-k1} with initial measures $\rho_0^N$ and
$\rho_0$ gives \eqref{eq:particle-tracking-k1}. The almost-sure convergence of
$W_1(\rho_0^N,\rho_0)$ for compactly supported i.i.d.\ samples, together with
the fact that $\theta_i^0\in B_0$ almost surely in that case, completes the
proof.
\end{proof}

\begin{remark}[Why the theorem is formulated on a common initialization box]
The confinement family $(B_t)_{t\in[0,T]}$ is tied to the common initialization
box $B_0$ extracted from $\rho_0$. This is exactly the level needed for the
two-run common-path application: both empirical initializations lie in $B_0$
almost surely, and both trained runs therefore evolve inside the same
deterministic confinement family.
\end{remark}

\begin{corollary}[Two-run common path]
\label{cor:complete-two-run-k1}
Assume Assumption~\ref{ass:complete-full-batch-k1}. Let
$\rho_{A,t}^N,\rho_{B,t}^N$ be two independent particle systems of the form
\eqref{eq:particle-system-k1}, with i.i.d.\ initial particles sampled from the
same law $\rho_0$. Then
\begin{equation}
  \sup_{t\in[0,T]}W_1(\rho_{A,t}^N,\rho_{B,t}^N)
  \le
  e^{2L_TT}
  \Big(
    W_1(\rho_{A,0}^N,\rho_0)+W_1(\rho_{B,0}^N,\rho_0)
  \Big),
  \label{eq:two-run-common-path-k1}
\end{equation}
and therefore
\[
  \sup_{t\in[0,T]}W_1(\rho_{A,t}^N,\rho_{B,t}^N)\to 0
\]
almost surely, hence also in probability, as $N\to\infty$.
\end{corollary}

\begin{proof}
Apply Theorem~\ref{thm:complete-full-batch-k1} to each run and use the
triangle inequality.
\end{proof}

\begin{remark}[Relation to the Ferbach training-side hypothesis]
\label{rem:ferbach-training-side}
Theorem~\ref{thm:complete-full-batch-k1} and
Corollary~\ref{cor:complete-two-run-k1} are the continuous-time full-batch
analogue of the training-side bridge used by Ferbach et al.~\cite{Ferbach2024}:
two wide runs started from the same law both track the same deterministic
trajectory on any fixed finite horizon. The present note proves this bridge for
the practical natural-softplus Gaussian head rather than assuming a scalar
convex output loss.
\end{remark}

\section{Deterministic Aligned Barrier Theory}

This section is the interpolation side. The only remaining gap between
parameter interpolation and convex interpolation in natural-output space is the
combination of vector-head transfer and softplus curvature.

\begin{lemma}[Automatic strip control along the interpolation path]
\label{lem:auto-strip-k1}
For all $x\in\mc{X}$ and all $t\in[0,1]$,
\[
  \norm{z_A(x)}_\infty,\ \norm{z_B(x)}_\infty,\ \norm{z_t(x)}_\infty,\
  \norm{\hat z_t(x)}_\infty
  \le
  K_\varphi M_V^\infty.
\]
Consequently,
\[
  \eta_{\Theta_A}(x),\ \eta_{\Theta_B}(x),\ \eta_{\Theta_t}(x),\
  \hat\eta_t(x),\ \tilde\eta_t(x)
  \in
  D_{K_\varphi M_V^\infty,\lambda_{\min}}
  \qquad
  \forall x,\ \forall t.
\]
\end{lemma}

\begin{proof}
For the first coordinate,
\[
  |u_{\Theta_A}(x)|
  \le
  \frac1N\sum_{i=1}^N |a_i^A|\,|\varphi(\xi_i^A,x)|
  \le
  K_\varphi
  \Bigl(\frac1N\sum_{i=1}^N (a_i^A)^2\Bigr)^{1/2}
  \le
  K_\varphi M_V^\infty.
\]
The same estimate holds for $s_{\Theta_A}$, and likewise for network $B$.
For the interpolated coefficients, convexity of the square gives
\[
  \frac1N\sum_{i=1}^N \bigl((1-t)a_i^A+ta_i^B\bigr)^2
  \le
  (1-t)\frac1N\sum_{i=1}^N (a_i^A)^2
  +
  t\frac1N\sum_{i=1}^N (a_i^B)^2
  \le
  (M_V^\infty)^2,
\]
and similarly for the $c$ coordinates. This proves the raw bound for $z_t$.
The same bound for $\hat z_t$ is immediate.

For the natural outputs, the first coordinate is unchanged and the second
coordinate is always $\le -\lambda_{\min}$ by
Proposition~\ref{prop:softplus-domain-k1}.
\end{proof}

\begin{proposition}[Softplus Jensen gap]
\label{prop:softplus-jensen-k1}
For any $s_A,s_B\in\R$ and $t\in[0,1]$,
\begin{equation}
  0
  \le
  (1-t)\psi_{\rm sp}(s_A)+t\psi_{\rm sp}(s_B)
  -
  \psi_{\rm sp}((1-t)s_A+ts_B)
  \le
  \frac18\,t(1-t)(s_A-s_B)^2.
  \label{eq:softplus-jensen-k1}
\end{equation}
Consequently, for every $x\in\mc{X}$,
\begin{equation}
  \norm{\tilde\eta_t(x)-\hat\eta_t(x)}_\infty
  \le
  \frac18\,t(1-t)\bigl(s_{\Theta_A}(x)-s_{\Theta_B}(x)\bigr)^2.
  \label{eq:natural-jensen-gap-k1}
\end{equation}
\end{proposition}

\begin{proof}
The function $\psi_{\rm sp}$ is convex and satisfies
\[
  \psi''_{\rm sp}(s)=\frac{e^s}{(1+e^s)^2}\le \frac14.
\]
For any $C^2$ convex function with second derivative bounded by $L$, the
interpolation error is at most $\frac{L}{2}t(1-t)(s_A-s_B)^2$. Applying this
with $L=\frac14$ gives \eqref{eq:softplus-jensen-k1}. Since only the second
natural coordinate is nonlinear, \eqref{eq:natural-jensen-gap-k1} is exactly
the same estimate written for $\tilde\eta_t$ and $\hat\eta_t$.
\end{proof}

\begin{proposition}[Direct vector-head transfer]
\label{prop:vector-head-transfer-k1}
For the aligned pair $(\Theta_A,\Theta_B)$,
\begin{equation}
  \sup_{t\in[0,1]}\ \sup_{x\in\mc{X}}
  \norm{z_t(x)-\hat z_t(x)}_\infty
  \le
  \frac12 K_{\nabla}M_V^\infty\sqrt{B_N}.
  \label{eq:vector-transfer-k1}
\end{equation}
In particular,
\[
  \sup_{t\in[0,1]}
  \Bigl(
    \E\big[\norm{z_t(X)-\hat z_t(X)}_\infty^2\big]
  \Bigr)^{1/2}
  \le
  \frac12 K_{\nabla}M_V^\infty\sqrt{B_N}.
\]
\end{proposition}

\begin{proof}
Set
\[
  \xi_i(t):=(1-t)\xi_i^A+t\xi_i^B,
  \qquad
  \Delta\xi_i:=\xi_i^A-\xi_i^B.
\]
Let $v_i^A=(a_i^A,c_i^A)$ and $v_i^B=(a_i^B,c_i^B)$. Then
\begin{align*}
  z_t(x)-\hat z_t(x)
  &=
  \frac1N\sum_{i=1}^N
  \Bigl[
    (1-t)v_i^A\bigl(\varphi(\xi_i(t),x)-\varphi(\xi_i^A,x)\bigr)
    \\
  &\qquad\qquad\qquad
    +
    tv_i^B\bigl(\varphi(\xi_i(t),x)-\varphi(\xi_i^B,x)\bigr)
  \Bigr].
\end{align*}
Since
\[
  \norm{\xi_i(t)-\xi_i^A}=t\norm{\Delta\xi_i},
  \qquad
  \norm{\xi_i(t)-\xi_i^B}=(1-t)\norm{\Delta\xi_i},
\]
the uniform gradient bound gives
\[
  |\varphi(\xi_i(t),x)-\varphi(\xi_i^A,x)|
  \le
  tK_{\nabla}\norm{\Delta\xi_i},
  \qquad
  |\varphi(\xi_i(t),x)-\varphi(\xi_i^B,x)|
  \le
  (1-t)K_{\nabla}\norm{\Delta\xi_i}.
\]
Taking the $j$th coordinate of the vector head and applying Cauchy--Schwarz,
\[
  |(z_t(x)-\hat z_t(x))_j|
  \le
  2t(1-t)K_{\nabla}M_V^\infty\sqrt{B_N}
  \le
  \frac12 K_{\nabla}M_V^\infty\sqrt{B_N}.
\]
Taking the maximum over $j\in\{1,2\}$ proves the claim.
\end{proof}

\begin{proposition}[Natural-output interpolation gap]
\label{prop:natural-output-gap-k1}
For every $t\in[0,1]$,
\begin{equation}
  \E\big[\norm{\eta_{\Theta_t}(X)-\hat\eta_t(X)}_\infty\big]
  \le
  \frac12 K_{\nabla}M_V^\infty\sqrt{B_N}
  +
  \frac18\,t(1-t)\Delta_{s,N}.
  \label{eq:natural-gap-k1}
\end{equation}
\end{proposition}

\begin{proof}
By the triangle inequality,
\[
  \norm{\eta_{\Theta_t}(x)-\hat\eta_t(x)}_\infty
  \le
  \norm{\eta_{\Theta_t}(x)-\tilde\eta_t(x)}_\infty
  +
  \norm{\tilde\eta_t(x)-\hat\eta_t(x)}_\infty.
\]
The map $\Gamma(u,s)=(u,-\psi_{\rm sp}(s))$ is $1$-Lipschitz in the
$\ell_\infty$ norm because its first coordinate is the identity and
$0<\psi'_{\rm sp}(s)<1$. Hence
\[
  \norm{\eta_{\Theta_t}(x)-\tilde\eta_t(x)}_\infty
  \le
  \norm{z_t(x)-\hat z_t(x)}_\infty,
\]
and Proposition~\ref{prop:vector-head-transfer-k1} bounds the expectation of
this term by $\frac12 K_{\nabla}M_V^\infty\sqrt{B_N}$. The second term is
controlled by Proposition~\ref{prop:softplus-jensen-k1}.
\end{proof}

\begin{theorem}[Deterministic aligned barrier theorem]
\label{thm:deterministic-lmc-k1}
For every fixed neuron permutation $\pi$,
\begin{equation}
  \cB(\Theta_A,\Theta_B;\pi)
  \le
  C_\eta(K_\varphi M_V^\infty,\lambda_{\min},Y_\ast)
  \left[
    \frac12 K_{\nabla}M_V^\infty\sqrt{B_N}
    +
    \frac1{32}\Delta_{s,N}
  \right].
  \label{eq:deterministic-barrier-k1}
\end{equation}
\end{theorem}

\begin{proof}
Fix $t\in[0,1]$. By Lemma~\ref{lem:auto-strip-k1}, all relevant outputs lie in
$D_{K_\varphi M_V^\infty,\lambda_{\min}}$, so
Proposition~\ref{prop:convex-lipschitz-strip-k1} applies there. First use the
Lipschitz bound:
\[
  \E[\ell_\eta(\eta_{\Theta_t}(X);Y)]
  \le
  \E[\ell_\eta(\hat\eta_t(X);Y)]
  +
  C_\eta(K_\varphi M_V^\infty,\lambda_{\min},Y_\ast)
  \E[\norm{\eta_{\Theta_t}(X)-\hat\eta_t(X)}_\infty].
\]
Now apply Proposition~\ref{prop:natural-output-gap-k1}. The remaining term is
handled by convexity:
\[
  \ell_\eta(\hat\eta_t(X);Y)
  \le
  (1-t)\ell_\eta(\eta_{\Theta_A}(X);Y)
  +
  t\ell_\eta(\eta_{\Theta_B}(X);Y).
\]
Taking expectations and then $\sup_{t\in[0,1]}$ gives
\eqref{eq:deterministic-barrier-k1}, since $\sup_t t(1-t)=1/4$.
\end{proof}

\begin{remark}[The vector-head issue]
\label{rem:vector-head-issue}
Even though $K=1$, the present head is not scalar. The pair $(u,s)$ is already
a two-coordinate output, so Ferbach's scalar-head transfer step is not
sufficient by itself. Proposition~\ref{prop:vector-head-transfer-k1} is the
appropriate replacement: it couples the whole raw two-coordinate head directly.
\end{remark}

\section{From the Common Path to Linear Mode Connectivity}

The training-side theorem gives two independent runs that are close in
Wasserstein distance. This section converts that common-path input into an
interpolation barrier bound. The argument itself is agnostic to whether the
common path came from the SGD theorem above or from the continuous-time
full-batch baseline.

\begin{lemma}[Permutation selected by empirical Wasserstein distance]
\label{lem:w1-permutation-k1}
Let
\[
  \rho_A^N:=\frac1N\sum_{i=1}^N\delta_{\theta_i^A},
  \qquad
  \rho_B^N:=\frac1N\sum_{i=1}^N\delta_{\theta_i^B},
  \qquad
  \theta_i^X:=(a_i^X,c_i^X,\xi_i^X).
\]
Let $\varepsilon_N:=W_1(\rho_A^N,\rho_B^N)$, where $W_1$ is computed with the
Euclidean cost on $\R^{2+p}$. Then there exists a permutation $\pi_N$ such
that
\begin{equation}
  \frac1N\sum_{i=1}^N\norm{\theta_i^A-\theta_{\pi_N(i)}^B}
  =
  \varepsilon_N.
  \label{eq:w1-permutation-k1}
\end{equation}
After relabeling network $B$ by this permutation, the alignment quantities
satisfy
\begin{align}
  B_N
  &\le
  2R_T\,\varepsilon_N,
  \label{eq:w1-to-BN-k1}
  \\
  \Delta_{s,N}
  &\le
  \bigl(K_\varphi+C_TK_{\nabla}\bigr)^2\varepsilon_N^2,
  \label{eq:w1-to-scale-k1}
\end{align}
where $A_T,C_T,R_T$ are the deterministic envelopes from
Proposition~\ref{prop:softplus-confinement-k1} and
$M_{V,T}:=\max\{A_T,C_T\}$.
\end{lemma}

\begin{proof}
Because both empirical measures place mass $1/N$ at each atom, the optimal
transport plan for $W_1$ is realized by a permutation, proving
\eqref{eq:w1-permutation-k1}.

For \eqref{eq:w1-to-BN-k1}, note that
\[
  \norm{\xi_i^A-\xi_i^B}\le 2R_T
  \qquad\text{and}\qquad
  \norm{\xi_i^A-\xi_i^B}\le \norm{\theta_i^A-\theta_i^B}.
\]
Hence
\[
  \norm{\xi_i^A-\xi_i^B}^2
  \le
  2R_T\norm{\theta_i^A-\theta_i^B}.
\]
Averaging proves \eqref{eq:w1-to-BN-k1}.

For \eqref{eq:w1-to-scale-k1}, write
\begin{align*}
  s_{\Theta_A}(x)-s_{\Theta_B}(x)
  &=
  \frac1N\sum_{i=1}^N
  (c_i^A-c_i^B)\varphi(\xi_i^A,x)
  \\
  &\qquad
  +
  \frac1N\sum_{i=1}^N
  c_i^B\bigl(\varphi(\xi_i^A,x)-\varphi(\xi_i^B,x)\bigr).
\end{align*}
Using $|\varphi|\le K_\varphi$, $|c_i^B|\le C_T$, and the gradient bound on
$\varphi$,
\[
  |s_{\Theta_A}(x)-s_{\Theta_B}(x)|
  \le
  \bigl(K_\varphi+C_TK_{\nabla}\bigr)\varepsilon_N.
\]
Squaring and taking expectation in $X$ yields \eqref{eq:w1-to-scale-k1}.
\end{proof}

\begin{theorem}[LMC from a common path]
\label{thm:lmc-from-common-path-k1}
Assume the trained endpoints satisfy the deterministic bounds
\[
  |a_i^A|,\ |a_i^B|\le A_T,
  \qquad
  |c_i^A|,\ |c_i^B|\le C_T,
  \qquad
  \norm{\xi_i^A},\ \norm{\xi_i^B}\le R_T
  \qquad
  \forall i=1,\dots,N,
\]
and let
\[
  \varepsilon_N:=W_1(\rho_A^N,\rho_B^N).
\]
Let $\pi_N$ be the permutation from Lemma~\ref{lem:w1-permutation-k1}. Then
\begin{equation}
  \cB(\Theta_A,\Theta_B;\pi_N)
  \le
  C_\eta(K_\varphi M_{V,T},\lambda_{\min},Y_\ast)
  \left[
    \frac12 K_{\nabla}M_{V,T}\sqrt{2R_T}\,\varepsilon_N^{1/2}
    +
    \frac{(K_\varphi+C_TK_{\nabla})^2}{32}\,\varepsilon_N^2
  \right],
  \label{eq:lmc-common-path-bound-k1}
\end{equation}
In particular, if $\varepsilon_N\to 0$, then
\[
  \cB(\Theta_A,\Theta_B;\pi_N)\to 0
\]
in the same mode of convergence.
\end{theorem}

\begin{proof}
The displayed deterministic bounds imply $M_V^\infty\le M_{V,T}$. Apply
Theorem~\ref{thm:deterministic-lmc-k1} after relabeling network $B$ by
$\pi_N$, and then use \eqref{eq:w1-to-BN-k1} and \eqref{eq:w1-to-scale-k1} to
replace $B_N$ and $\Delta_{s,N}$ by Wasserstein-based quantities.
\end{proof}

\begin{corollary}[Finite-horizon SGD LMC from a shared initialization law]
\label{cor:sgd-lmc-k1}
Assume Assumption~\ref{ass:softplus-sgd-k1}. Let
$\hat\rho_{A,k}^N,\hat\rho_{B,k}^N$ be two independent width-$N$ SGD runs of
\eqref{eq:sgd-update-k1}, both initialized i.i.d.\ from the same law $\rho_0$
and trained over the same horizon $[0,T]$. Let $\pi_{N,k}$ be the permutation
from Lemma~\ref{lem:w1-permutation-k1} applied at step $k$, and let
$\Theta_{A,k},\Theta_{B,k}$ denote the corresponding parameter vectors. Then,
with high probability and uniformly for all $k\le T/\varepsilon$,
\begin{align*}
  \cB(\Theta_{A,k},\Theta_{B,k};\pi_{N,k})
  &\le
  C_\eta(K_\varphi M_{V,T}^{\rm sgd},\lambda_{\min},Y_\ast)
  \Biggl[
    \frac12 K_{\nabla}M_{V,T}^{\rm sgd}\sqrt{2R_T^{\rm sgd}}\,
    \bigl(2\delta_{N,\varepsilon}(T)\bigr)^{1/2}
    \\
  &\qquad\qquad
    +
    \frac{(K_\varphi+C_T^{\rm sgd}K_{\nabla})^2}{8}\,
    \delta_{N,\varepsilon}(T)^2
  \Biggr].
\end{align*}
where
\[
  M_{V,T}^{\rm sgd}:=\max\{A_T^{\rm sgd},C_T^{\rm sgd}\}.
\]
In particular, for every fixed finite horizon, the aligned interpolation
barrier vanishes in probability as $N\to\infty$ and $\varepsilon\downarrow 0$.
\end{corollary}

\begin{proof}
Apply Corollary~\ref{cor:sgd-common-path-k1} and then
Theorem~\ref{thm:lmc-from-common-path-k1}, using the deterministic confinement
bounds from Proposition~\ref{prop:sgd-confinement-k1}.
\end{proof}

\begin{corollary}[Continuous-time full-batch specialization]
\label{cor:full-batch-lmc-k1}
Assume Assumption~\ref{ass:complete-full-batch-k1}. Let
$\rho_{A,t}^N,\rho_{B,t}^N$ be two independent full-batch runs started from the
same initialization law $\rho_0$. Then, for every fixed $t\in[0,T]$, the
empirical measures satisfy
\[
  W_1(\rho_{A,t}^N,\rho_{B,t}^N)\to 0
\]
almost surely, and therefore the aligned interpolation barrier at time $t$
vanishes by Theorem~\ref{thm:lmc-from-common-path-k1}.
\end{corollary}

\begin{proof}
Apply Corollary~\ref{cor:complete-two-run-k1} and then
Theorem~\ref{thm:lmc-from-common-path-k1}.
\end{proof}

\begin{remark}[Barrier decomposition]
\label{rem:barrier-decomposition-k1}
The bound \eqref{eq:lmc-common-path-bound-k1} splits into
\begin{enumerate}[label=(\alph*)]
  \item the transfer term, of order $\varepsilon_N^{1/2}$;
  \item the softplus Jensen term, of order $\varepsilon_N^2$.
\end{enumerate}
Thus, in the first $K=1$ theorem, the rate-limiting step is the
parameter-to-output transfer rather than the natural-output geometry.
\end{remark}

\section{Rates}

The preceding theorem is stated in terms of $\varepsilon_N=W_1(\rho_A^N,\rho_B^N)$.
To recover the more explicit width rates familiar from Ferbach-style arguments,
one inserts sharper alignment information.

\begin{proposition}[From hidden-law concentration to $B_N=O(d/N)$]
\label{prop:BN-from-w2-k1}
Let
\[
  \hat\nu_A^N:=\frac1N\sum_{i=1}^N \delta_{\xi_i^A},
  \qquad
  \hat\nu_B^N:=\frac1N\sum_{i=1}^N \delta_{\xi_i^B},
\]
and let $\mu_{\xi,\ast}$ be a reference law on $\R^d$. Suppose that for some
deterministic sequence $\eta_N\downarrow 0$,
\[
  W_2(\hat\nu_A^N,\mu_{\xi,\ast})\le \eta_N,
  \qquad
  W_2(\hat\nu_B^N,\mu_{\xi,\ast})\le \eta_N.
\]
If the neuron permutation $\pi_N$ is chosen by hidden-weight optimal transport,
then
\begin{equation}
  B_N = W_2^2(\hat\nu_A^N,\hat\nu_B^N)\le 4\eta_N^2.
  \label{eq:BN-via-w2-k1}
\end{equation}
In particular, if
\[
  \eta_N=O(\kappa_\xi\sqrt{d/N}),
\]
then
\[
  B_N=O(\kappa_\xi^2 d/N).
\]
\end{proposition}

\begin{proof}
Because the hidden empirical laws have equal weights $1/N$, the $W_2$-optimal
coupling between them is realized by a permutation, and the corresponding
squared cost is exactly $B_N$. The triangle inequality gives
\[
  W_2(\hat\nu_A^N,\hat\nu_B^N)
  \le
  W_2(\hat\nu_A^N,\mu_{\xi,\ast}) + W_2(\hat\nu_B^N,\mu_{\xi,\ast})
  \le
  2\eta_N.
\]
Squaring proves \eqref{eq:BN-via-w2-k1}.
\end{proof}

\begin{corollary}[Baseline rate]
\label{cor:baseline-rate-k1}
Suppose that along a width sequence,
\[
  B_N=O(d/N),
  \qquad
  \Delta_{s,N}=O(1/N).
\]
Then Theorem~\ref{thm:deterministic-lmc-k1} yields
\[
  \cB(\Theta_A,\Theta_B;\pi)
  =
  O\!\left(\sqrt{\frac{d}{N}}+\frac1N\right).
\]
In particular, for fixed input dimension $d$,
\[
  \cB(\Theta_A,\Theta_B;\pi)=O(N^{-1/2}).
\]
\end{corollary}

\begin{proof}
Substitute the displayed rates into \eqref{eq:deterministic-barrier-k1}.
\end{proof}

\begin{remark}[Ferbach-style specialization]
\label{rem:ferbach-style-rate-k1}
Proposition~\ref{prop:BN-from-w2-k1} is the precise point where the present
note meets the Ferbach-style alignment theory. If the trained hidden empirical
laws obey a concentration bound
\[
  \eta_N=O(\kappa_\xi\sqrt{d/N}),
\]
then hidden-weight optimal transport gives
\[
  B_N=O(\kappa_\xi^2 d/N).
\]
Substituting this into Corollary~\ref{cor:baseline-rate-k1} yields the
explicit baseline barrier rate
\[
  \cB(\Theta_A,\Theta_B;\pi_N)
  =
  O\!\left(\kappa_\xi\sqrt{\frac{d}{N}}+\frac1N\right).
\]
\end{remark}

The older LMC draft also identified a faster route: if the trained head varies
regularly with the hidden feature, the transfer term improves from
$O(\sqrt{B_N})$ to $O(B_N)$.

\begin{definition}[Per-coordinate head mismatch]
\label{def:head-mismatch-k1}
Define
\[
  M_{\Delta V}^\infty
  :=
  \max_{j\in\{1,2\}}
  \Bigl(
    \frac1N\sum_{i=1}^N \bigl((v_i^A-v_i^B)_j\bigr)^2
  \Bigr)^{1/2},
  \qquad
  v_i^X:=(a_i^X,c_i^X)\in\R^2.
\]
\end{definition}

\begin{assumption}[Conditional head regularity]
\label{ass:head-regularity-k1}
There exist a map $g:\R^d\to\R^2$ and residuals $r_i^A,r_i^B\in\R^2$ such that
\[
  v_i^A=g(\xi_i^A)+r_i^A,
  \qquad
  v_i^B=g(\xi_i^B)+r_i^B,
\]
and
\[
  \norm{g(\xi)-g(\bar\xi)}_\infty
  \le
  L_g\norm{\xi-\bar\xi}
  \qquad
  \forall \xi,\bar\xi\in\R^d.
\]
Define the residual mismatch
\[
  \mc{R}_N
  :=
  \max_{j\in\{1,2\}}
  \Bigl(
    \frac1N\sum_{i=1}^N \bigl((r_i^A-r_i^B)_j\bigr)^2
  \Bigr)^{1/2}.
\]
\end{assumption}

\begin{lemma}[Head mismatch from conditional regularity]
\label{lem:head-mismatch-regular-k1}
Under Assumption~\ref{ass:head-regularity-k1},
\begin{equation}
  M_{\Delta V}^\infty
  \le
  L_g\sqrt{B_N}+\mc{R}_N.
  \label{eq:head-mismatch-regular-k1}
\end{equation}
\end{lemma}

\begin{proof}
For each coordinate $j\in\{1,2\}$,
\[
  |(v_i^A-v_i^B)_j|
  \le
  |(g(\xi_i^A)-g(\xi_i^B))_j|
  +
  |(r_i^A-r_i^B)_j|
  \le
  L_g\norm{\xi_i^A-\xi_i^B}+|(r_i^A-r_i^B)_j|.
\]
Take the empirical $\ell_2$ norm in $i$ and the maximum over $j$.
\end{proof}

\begin{proposition}[Refined vector-head transfer]
\label{prop:refined-transfer-k1}
Assume Assumption~\ref{ass:complete-full-batch-k1}. Assume in addition that
$\nabla_\xi\varphi$ is Lipschitz in $\xi$ with constant $L_{\nabla}$
uniformly in $x$. Then
\begin{equation}
  \sup_{t\in[0,1]}\ \sup_{x\in\mc{X}}
  \norm{z_t(x)-\hat z_t(x)}_\infty
  \le
  \frac{K_{\nabla}}{4}M_{\Delta V}^\infty\sqrt{B_N}
  +
  \frac{L_{\nabla}}{8}M_{V,T}B_N,
  \label{eq:refined-transfer-k1}
\end{equation}
where $M_{V,T}:=\max\{A_T,C_T\}$.
\end{proposition}

\begin{proof}
Write
\[
  \phi_i^A(x):=\varphi(\xi_i^A,x),
  \qquad
  \phi_i^B(x):=\varphi(\xi_i^B,x),
  \qquad
  \phi_i^t(x):=\varphi((1-t)\xi_i^A+t\xi_i^B,x).
\]
Then
\begin{align*}
  z_t(x)-\hat z_t(x)
  &=
  \frac1N\sum_{i=1}^N
  \Bigl(
    v_i^t\phi_i^t-(1-t)v_i^A\phi_i^A-tv_i^B\phi_i^B
  \Bigr)
  \\
  &=
  \frac1N\sum_{i=1}^N
  v_i^t\bigl(\phi_i^t-(1-t)\phi_i^A-t\phi_i^B\bigr)
  \\
  &\qquad
  +
  t(1-t)\frac1N\sum_{i=1}^N
  (v_i^B-v_i^A)(\phi_i^A-\phi_i^B),
\end{align*}
where $v_i^t:=(1-t)v_i^A+tv_i^B$.

For the second sum, use
\[
  |\phi_i^A(x)-\phi_i^B(x)|
  \le
  K_{\nabla}\norm{\xi_i^A-\xi_i^B}
\]
and Cauchy--Schwarz to obtain
\[
  t(1-t)\frac1N\sum_{i=1}^N
  |(v_i^B-v_i^A)_j|\,|\phi_i^A(x)-\phi_i^B(x)|
  \le
  \frac{K_{\nabla}}{4}M_{\Delta V}^\infty\sqrt{B_N}.
\]

For the first sum, the Lipschitz continuity of $\nabla_\xi\varphi$ gives the
second-order interpolation remainder
\[
  \bigl|
    \phi_i^t(x)-(1-t)\phi_i^A(x)-t\phi_i^B(x)
  \bigr|
  \le
  \frac{L_{\nabla}}{2}\,t(1-t)\norm{\xi_i^A-\xi_i^B}^2.
\]
Since $|(v_i^t)_j|\le M_{V,T}$ under the finite-horizon confinement bound,
\[
  \frac1N\sum_{i=1}^N
  |(v_i^t)_j|
  \bigl|
    \phi_i^t(x)-(1-t)\phi_i^A(x)-t\phi_i^B(x)
  \bigr|
  \le
  \frac{L_{\nabla}}{8}M_{V,T}B_N.
\]
Adding the two estimates and taking the maximum over $j$ proves the claim.
\end{proof}

\begin{proposition}[Endpoint raw-scale rate]
\label{prop:scale-rate-k1}
Assume Assumption~\ref{ass:complete-full-batch-k1}. Then
\begin{equation}
  \Delta_{s,N}^{1/2}
  \le
  K_\varphi M_{\Delta V}^\infty + C_TK_{\nabla}\sqrt{B_N}.
  \label{eq:scale-rate-k1}
\end{equation}
Consequently, if
\[
  M_{\Delta V}^\infty=O(\sqrt{B_N}),
\]
then
\[
  \Delta_{s,N}=O(B_N).
\]
\end{proposition}

\begin{proof}
Exactly as in the proof of Lemma~\ref{lem:w1-permutation-k1},
\begin{align*}
  s_{\Theta_A}(x)-s_{\Theta_B}(x)
  &=
  \frac1N\sum_{i=1}^N
  (c_i^A-c_i^B)\varphi(\xi_i^A,x)
  \\
  &\qquad
  +
  \frac1N\sum_{i=1}^N
  c_i^B\bigl(\varphi(\xi_i^A,x)-\varphi(\xi_i^B,x)\bigr).
\end{align*}
The first term is bounded by $K_\varphi M_{\Delta V}^\infty$ and the second by
$C_TK_{\nabla}\sqrt{B_N}$. Take the supremum in $x$ and then the
$L^2(P_X)$ norm.
\end{proof}

\begin{corollary}[Fast rate under conditional head regularity]
\label{cor:fast-rate-k1}
Assume Assumption~\ref{ass:complete-full-batch-k1} and
Assumption~\ref{ass:head-regularity-k1}. Assume also that
$\nabla_\xi\varphi$ is Lipschitz with constant $L_{\nabla}$ and that
\[
  \mc{R}_N=O(\sqrt{B_N}).
\]
Then
\[
  \cB(\Theta_A,\Theta_B;\pi)=O(B_N).
\]
In particular, if
\[
  B_N=O(d/N),
\]
then
\[
  \cB(\Theta_A,\Theta_B;\pi)=O(d/N),
\]
which is $O(1/N)$ for fixed input dimension $d$.
\end{corollary}

\begin{proof}
Lemma~\ref{lem:head-mismatch-regular-k1} implies
$M_{\Delta V}^\infty=O(\sqrt{B_N})$. Proposition~\ref{prop:refined-transfer-k1}
then yields
\[
  \sup_t\sup_x \norm{z_t(x)-\hat z_t(x)}_\infty = O(B_N).
\]
Because $\Gamma$ is $1$-Lipschitz, the same bound holds for
$\eta_{\Theta_t}-\tilde\eta_t$. Proposition~\ref{prop:scale-rate-k1} gives
$\Delta_{s,N}=O(B_N)$, hence Proposition~\ref{prop:softplus-jensen-k1} gives
\[
  \sup_t \E\big[\norm{\tilde\eta_t(X)-\hat\eta_t(X)}_\infty\big]=O(B_N).
\]
Repeating the proof of Theorem~\ref{thm:deterministic-lmc-k1} with these
refined estimates gives $\cB=O(B_N)$.
\end{proof}

\section{Conclusion}

For the single-Gaussian head, the mean-field-to-LMC program can be completed in
one clean argument.
\begin{enumerate}[label=(\roman*)]
  \item The practical natural-softplus head admits a finite-horizon SGD
  common-path theorem, with a complete continuous-time full-batch theorem as a
  fully proved baseline.
  \item The interpolation barrier is controlled by exact convexity in natural
  coordinates together with two explicitly quantified defects: vector-head
  transfer and softplus curvature.
  \item Combining these two facts yields a Ferbach-style LMC theorem for the
  $K=1$ distributional model.
\end{enumerate}

This paper therefore provides the clean baseline that is missing from the
broader mixture setting: it isolates what survives from the convex-loss LMC
theory once one replaces the scalar head by a practice-faithful distributional
head, while still avoiding the additional gauge and non-identifiability issues
that appear for $K>1$.

\begin{thebibliography}{99}

\bibitem{Ferbach2024}
Damien Ferbach, Baptiste Goujaud, Gauthier Gidel, and Aymeric Dieuleveut.
\newblock Proving linear mode connectivity of neural networks via optimal transport.
\newblock In \emph{Proceedings of The 27th International Conference on Artificial Intelligence and Statistics (AISTATS)}, volume 238 of \emph{Proceedings of Machine Learning Research}, pages 3853--3861, 2024.

\bibitem{Frankle2020}
Jonathan Frankle, Gintare Karolina Dziugaite, Daniel Roy, and Michael Carbin.
\newblock Linear mode connectivity and the lottery ticket hypothesis.
\newblock In \emph{Proceedings of the 37th International Conference on Machine Learning (ICML)}, pages 3259--3269, 2020.

\bibitem{Garipov2018}
Timur Garipov, Pavel Izmailov, Dmitrii Podoprikhin, Dmitry Vetrov, and Andrew Gordon Wilson.
\newblock Loss surfaces, mode connectivity, and fast ensembling of {DNN}s.
\newblock In \emph{Advances in Neural Information Processing Systems 31 (NeurIPS)}, pages 8789--8798, 2018.

\bibitem{MeiMisiakiewiczMontanari2019}
Song Mei, Theodor Misiakiewicz, and Andrea Montanari.
\newblock Mean-field theory of two-layer neural networks: dimension-free bounds and kernel limit.
\newblock In \emph{Proceedings of the 32nd Conference on Learning Theory (COLT)}, volume 99 of \emph{Proceedings of Machine Learning Research}, pages 2388--2464, 2019.

\bibitem{SirignanoSpiliopoulos2020}
Justin Sirignano and Konstantinos Spiliopoulos.
\newblock Mean field analysis of neural networks: a law of large numbers.
\newblock \emph{SIAM Journal on Applied Mathematics}, 80(2):725--752, 2020.

\bibitem{WainwrightJordan2008}
Martin J. Wainwright and Michael I. Jordan.
\newblock Graphical models, exponential families, and variational inference.
\newblock \emph{Foundations and Trends in Machine Learning}, 1(1--2):1--305, 2008.

\end{thebibliography}

\end{document}
